\subsection{Anti-gaming: serving-unavailable features and the leak we found}
\label{sec:q9}

\eb{} ships three lifetime totals per article (views,
page views, read time) summed over the whole collection period, so for any impression they include
the future and cannot exist when serving. We trained the same ranker on the same rows with and
without them:

\begin{table*}[t]
\centering\small
\begin{tabular}{@{}lc@{}}
\toprule
\eb{} small, val, 64,365 impressions & per-impression AUC \\
\midrule
causal features only (servable) & \QnineCausal \\
plus lifetime totals (not servable) & \QnineLeaky \\
\midrule
difference, paired bootstrap & \QnineDelta{} \QnineCI \\
\bottomrule
\end{tabular}
\caption{Metrics with and without serving-unavailable features, measured on the 10-feature
model with the 6\,h window.}
\label{tab:q9}
\end{table*}

In Assignment~1 the same totals were worth $+$0.042 to $+$0.075, more than our whole semantic axis,
because they competed against systems with no popularity signal at all. Here they compete with
click counts restricted to the past, which capture most of the same information legitimately.
Cheating would buy \QnineDelta{}; the behavioural axis earns \EBGainTest{} legitimately, roughly an
order of magnitude more. Both arms are the 10-feature model, so this prices the leak against that
round rather than the shipped 22. End to end the servable re-ranker (\EBFinalTest{}) far outscores
stage one with lifetime popularity added (\EBLeakyTest{}). \mind{} cannot run this comparison at
all: its lifetime columns are null for all 65,238 articles.

Two tests guard the boundary: one asserts that no click, exposure or
engagement value used for an impression comes from at or after its timestamp, on every split; the
other fails if any scoring module reads a lifetime column. Each was verified by injecting a
violation and watching it fail, after we found checks that passed without checking anything
(Section~\ref{sec:lessons}).

The first recounts popularity against the 6\,h window the shipped feature actually uses, with a \emph{two-sided} mask and a different code path from the builder's binary search. Two-sided
is load-bearing: asserting only ``no more than the clicks before $t$'' is satisfied by a window
reaching past $t$ so long as it stays narrow, so the weaker check passes on the leak it exists to
catch. We confirmed the stronger one is live by mutation rather than argument, since changing the
window to 12\,h or unbounded without rebuilding makes both boundary tests fail. The suite is 29 passed, 6 skipped, 0 failed.

We also found a real leak in our own code, and fixing it raised accuracy. Merging \eb{}'s two
history blocks described 99.8\% of validation and 99.9\% of training impressions with history from
their own future; test was clean at 0.0\%. It was worth $+$0.0967 on the affected
feature, and after the fix test AUC \emph{rose} from 0.7429 to 0.7461, because test features were
always clean and only the model changed. How we found it, and the evidence that the fix was right,
are in Appendix~\ref{app:leak}.
